\chapter{Wave Functions, Density Matrices, and Locality}

\lectureoverview{Bell's theorem is postponed for one essential reason: before asking what Bell forbids, we must say precisely what entanglement is. A bipartite state is entangled when it cannot be factored into one state for Alice and one for Bob. Wave functions make that criterion concrete, correlations reveal its consequences, and the singlet supplies the decisive example. The reduced density matrix then becomes the complete local description of one part of an entangled whole. With that machinery in hand, measurement appears as an entangling interaction, the system--observer boundary can be moved without changing predictions, and the apparent tension between entanglement and locality can be stated sharply: distant choices change conditional correlations, but not the unconditioned statistics available for signaling.}

\section{What Entanglement Actually Means}

It is tempting to proceed directly to Bell's theorem. The theorem is beautiful, but its significance cannot be separated from the structure it tests. The more urgent question is therefore not yet whether quantum mechanics violates a Bell inequality, but what an entangled state says about a composite system.

A common phrase is that two entangled systems \emph{share a state}. That phrase does not distinguish anything. Every composite system has a state, whether entangled or not. The distinction is factorization. If \(A\) and \(B\) have Hilbert spaces \(\mathcal H_A\) and \(\mathcal H_B\), the state of the pair lies in
\begin{equation}
\mathcal H_{AB}=\mathcal H_A\otimes\mathcal H_B.
\end{equation}
A pure state is a product state when it can be written
\begin{equation}
\ket{\Psi}_{AB}=\ket{\psi}_A\otimes\ket{\phi}_B.
\label{eq:l7-product-state}
\end{equation}
If no such pair of factors exists, \(\ket{\Psi}_{AB}\) is entangled.

\begin{classroomqa}
\audiencequestion{Does entanglement simply mean that both systems occur in one state vector?}
\lecturerresponse{No. A product state also belongs to the Hilbert space of the whole pair. Entanglement is the stronger, basis-independent statement that the vector of the whole cannot be written as one vector for Alice times one vector for Bob.}
\end{classroomqa}

\begin{classroomqa}
\audiencequestion{Could a clever change of basis turn an entangled state into a product state?}
\lecturerresponse{Not if the change of basis is local, meaning one basis change for Alice and another for Bob. Local basis changes alter the coordinates of each factor but do not alter whether the state factors. A completely arbitrary basis of the joint Hilbert space can hide the tensor-product structure, but it does not change which physical degrees of freedom belong to Alice and Bob.}
\end{classroomqa}

\section{Wave Functions as Coordinates in Hilbert Space}

The word \emph{wave function} carries historical imagery that is not part of its mathematical definition. A wave function is the coordinate representation of a state vector in a chosen basis. Let a complete set of commuting observables have simultaneous eigenvectors
\begin{equation}
\ket{a,b,c,\ldots}.
\end{equation}
The labels must be complete: taken together, they distinguish every vector in the basis. Then
\begin{align}
\ket{\Psi}
&=\sum_{a,b,c,\ldots}
\ket{a,b,c,\ldots}
\braket{a,b,c,\ldots}{\Psi} \\
&=\sum_{a,b,c,\ldots}
\psi(a,b,c,\ldots)\ket{a,b,c,\ldots},
\label{eq:l7-wavefunction-expansion}
\end{align}
where
\begin{equation}
\psi(a,b,c,\ldots)
=\braket{a,b,c,\ldots}{\Psi}.
\label{eq:l7-wavefunction-definition}
\end{equation}
For continuous labels the corresponding sums become integrals.

\begin{figure}[htbp]
\centering
\includegraphics[width=0.88\textwidth]{lecture_07_figure_01.png}
\caption{The wave function as the set of coefficients of a state in a complete basis.}
\lectureframe{00:07:00}
\end{figure}

For a particle in three-dimensional space, the commuting labels may be \(x,y,z\), and
\begin{equation}
\psi(x,y,z)=\braket{x,y,z}{\Psi}
\end{equation}
is Schr\"odinger's familiar position-space wave function. For one spin in the \(\sigma_z\) basis, the wave function is only the pair
\begin{equation}
\psi(u)=\braket{u}{\Psi},
\qquad
\psi(d)=\braket{d}{\Psi}.
\end{equation}
For two spins, two labels are needed:
\begin{equation}
\Psi(a,b)=\braket{a,b}{\Psi},
\qquad
a,b\in\{u,d\}.
\label{eq:l7-two-spin-wavefunction}
\end{equation}
Nothing in this definition requires a literal wave in space.

\begin{classroomqa}
\audiencequestion{Is changing the wave-function basis a Fourier transform?}
\lecturerresponse{A Fourier transform is the particular basis change between position and momentum representations. More generally, every change of orthonormal basis is a unitary transformation of the coordinates. Passing from the \(\sigma_z\) basis to the \(\sigma_x\) basis is analogous to a Fourier transform, though it is a finite-dimensional unitary matrix rather than the position--momentum integral transform.}
\end{classroomqa}

\begin{classroomqa}
\audiencequestion{If the wave function changes under a basis change, has the physical state changed?}
\lecturerresponse{No. The ket \(\ket{\Psi}\) is the abstract state. The numbers \(\psi(a,b,c,\ldots)\) are its coordinates in a selected basis. Changing coordinates changes those numbers while leaving the vector and all physical predictions unchanged.}
\end{classroomqa}

\section{Expectation Values in Components}

Before introducing density matrices, expose the component structure of an expectation value. Let \(L\) be an observable, with matrix elements
\begin{equation}
L_{a'a}=\bra{a'}L\ket{a}.
\label{eq:l7-matrix-elements}
\end{equation}
If \(\ket{\Psi}=\sum_a\psi(a)\ket{a}\), then
\begin{align}
\expect{L}
&=\bra{\Psi}L\ket{\Psi} \\
&=\sum_{a',a}
\psi^*(a')\bra{a'}L\ket{a}\psi(a) \\
&=\sum_{a',a}\psi^*(a')L_{a'a}\psi(a).
\label{eq:l7-expectation-components}
\end{align}

\begin{figure}[htbp]
\centering
\includegraphics[width=0.88\textwidth]{lecture_07_figure_02.png}
\caption{The expectation-value sandwich expanded into wave-function coefficients and matrix elements.}
\lectureframe{00:14:08}
\end{figure}

Now divide the system into Alice's part and Bob's part. Let \(a\) be a complete basis label for Alice and \(b\) one for Bob:
\begin{equation}
\ket{\Psi}_{AB}
=\sum_{a,b}\Psi(a,b)\ket{a,b}.
\label{eq:l7-composite-expansion}
\end{equation}
An observable belonging only to Alice is represented on the composite space by
\begin{equation}
L_A=L\otimes I_B,
\end{equation}
so its matrix elements are
\begin{equation}
\bra{a',b'}L_A\ket{a,b}
=L_{a'a}\delta_{b'b}.
\label{eq:l7-alice-local-matrix}
\end{equation}
Likewise an observable local to Bob has the form
\begin{equation}
M_B=I_A\otimes M,
\qquad
\bra{a',b'}M_B\ket{a,b}
=\delta_{a'a}M_{b'b}.
\label{eq:l7-bob-local-matrix}
\end{equation}
The identity in the untouched factor is the mathematical statement that Alice's operation acts on Alice and Bob's operation acts on Bob.

Substituting Equation~\eqref{eq:l7-alice-local-matrix} into the expectation value gives
\begin{equation}
\expect{L_A}
=\sum_{a',a,b}
\Psi^*(a',b)L_{a'a}\Psi(a,b).
\label{eq:l7-local-expectation}
\end{equation}
This formula already contains the reduced density matrix; we will reorganize it after first seeing what factorization does.

\begin{classroomqa}
\audiencequestion{Why is there only one Bob index left in Equation~\eqref{eq:l7-local-expectation}?}
\lecturerresponse{The identity on Bob's space contributes \(\delta_{b'b}\). It sets the bra and ket labels on Bob's side equal, so the double sum over \(b',b\) collapses to one sum over \(b\).}
\end{classroomqa}

\begin{classroomqa}
\audiencequestion{Does Alice need to know which basis Bob uses?}
\lecturerresponse{No physical prediction depends on Bob's basis. The sum over a complete Bob basis is a resolution of the identity. A different Bob basis changes the individual components \(\Psi(a,b)\), but not Alice's expectation value.}
\end{classroomqa}

\section{Product States and Correlations}

In a product state the joint wave function factorizes:
\begin{equation}
\Psi(a,b)=\psi_A(a)\phi_B(b).
\label{eq:l7-wavefunction-factorization}
\end{equation}
Insert this into Equation~\eqref{eq:l7-local-expectation}:
\begin{align}
\expect{L_A}
&=\sum_{a',a,b}
\psi_A^*(a')\phi_B^*(b)
L_{a'a}
\psi_A(a)\phi_B(b) \\
&=\left[
\sum_{a',a}\psi_A^*(a')L_{a'a}\psi_A(a)
\right]
\left[
\sum_b|\phi_B(b)|^2
\right] \\
&=\bra{\psi_A}L\ket{\psi_A}.
\label{eq:l7-product-local-expectation}
\end{align}
Bob's normalization makes his factor equal to one. Alice therefore has exactly the same local statistics she would have had if Bob's system had never been introduced.

\begin{figure}[htbp]
\centering
\includegraphics[width=0.88\textwidth]{lecture_07_figure_03.png}
\caption{Factorization of the joint wave function and the corresponding separation of local expectation values.}
\lectureframe{00:29:18}
\end{figure}

The product of one Alice observable and one Bob observable has matrix elements
\begin{equation}
\bra{a',b'}L_A M_B\ket{a,b}
=L_{a'a}M_{b'b}.
\end{equation}
For a product state its expectation value separates:
\begin{equation}
\expect{L_A M_B}
=\expect{L_A}\expect{M_B}.
\label{eq:l7-product-correlation-factorization}
\end{equation}
It is therefore natural to define the connected correlation
\begin{equation}
C(L,M)
=\expect{L_A M_B}
-\expect{L_A}\expect{M_B}.
\label{eq:l7-connected-correlation}
\end{equation}
Every product state has \(C(L,M)=0\) for every pair of local observables. Conversely, for a pure bipartite state, entanglement guarantees that some pair of local observables has a nonzero connected correlation.

\begin{classroomqa}
\audiencequestion{Does \(L_A M_B\) mean that Bob is measured first and Alice second?}
\lecturerresponse{No. Operators act on state vectors; apparatuses perform measurements. The product means that \(M_B\) acts and then \(L_A\) acts. Since the operators occupy different tensor factors, \([L_A,M_B]=0\), so their mathematical order is immaterial. A laboratory protocol for estimating the joint expectation is a separate matter.}
\end{classroomqa}

\begin{classroomqa}
\audiencequestion{Can the state be entangled even when one chosen correlation vanishes?}
\lecturerresponse{Certainly. One vanishing \(C(L,M)\) says only that this particular pair of observables is uncorrelated. To establish factorization for a pure state, all local cross-correlations must factorize. Partial correlation also need not mean partial entanglement in any simple one-number sense; the full reduced-density spectrum carries that information.}
\end{classroomqa}

\begin{classroomqa}
\audiencequestion{Is every correlated state entangled?}
\lecturerresponse{For the pure states under discussion, nonfactorizing correlations reveal entanglement. For mixed states, ordinary classical correlations can also make \(C(L,M)\neq0\). The later density-matrix criterion must therefore be qualified by whether the state of the whole is known to be pure.}
\end{classroomqa}

\section{The Singlet: A Definite Whole with Indefinite Parts}

Let \(\boldsymbol{\sigma}\) act on Bob's spin and \(\boldsymbol{\tau}\) on Alice's. Their singlet state is
\begin{equation}
\ket{S}
=\frac{1}{\sqrt2}
\left(\ket{u d}-\ket{d u}\right),
\label{eq:l7-singlet}
\end{equation}
where the first slot is Bob's and the second is Alice's. Neither spin has a definite \(z\)-component:
\begin{equation}
\expect{\sigma_z}=0,
\qquad
\expect{\tau_z}=0.
\label{eq:l7-singlet-local-z}
\end{equation}
Yet their product is certain:
\begin{equation}
\expect{\sigma_z\tau_z}=-1.
\label{eq:l7-singlet-product-z}
\end{equation}
Every \(z\)-measurement produces opposite results.

Rotational invariance makes the statement independent of direction. For any unit vector \(\mathbf n\),
\begin{equation}
\expect{\boldsymbol{\sigma}\cdot\mathbf n}
=\expect{\boldsymbol{\tau}\cdot\mathbf n}=0,
\qquad
\expect{
(\boldsymbol{\sigma}\cdot\mathbf n)
(\boldsymbol{\tau}\cdot\mathbf n)
}=-1.
\label{eq:l7-singlet-any-axis}
\end{equation}
More generally,
\begin{equation}
\expect{\sigma_i\tau_j}=-\delta_{ij}.
\label{eq:l7-singlet-correlation-tensor}
\end{equation}
Each part is completely unpolarized, while equal-axis outcomes are perfectly anticorrelated.

\begin{figure}[htbp]
\centering
\includegraphics[width=0.88\textwidth]{lecture_07_figure_04.png}
\caption{For the singlet, each local spin average vanishes while equal components have product \(-1\).}
\lectureframe{00:38:28}
\end{figure}

This is the characteristic reversal produced by entanglement. For an isolated spin in a pure state there is some direction along which its spin is definite. In the singlet, no direction makes either local spin definite. Nevertheless the composite state is pure and completely specified. We know as much as quantum mechanics permits about the whole and still cannot assign a pure state to either part.

\begin{classroomqa}
\audiencequestion{If the state of the whole is known exactly, why can we not infer a definite state for each spin?}
\lecturerresponse{Because the singlet is not a product. Its coefficients do not split into a Bob amplitude times an Alice amplitude. The information resides in relations between the parts rather than in two independently existing local wave functions.}
\end{classroomqa}

\begin{classroomqa}
\audiencequestion{Is an operator such as \(\sigma_z\) itself a possible state of the spin?}
\lecturerresponse{No. A state is a vector or, more generally, a density operator; \(\sigma_z\) is an observable acting on states. An eigenvector of \(\sigma_z\) is a state with a definite \(z\)-component, but the observable and its eigenvectors are different kinds of objects.}
\end{classroomqa}

\begin{classroomqa}
\audiencequestion{When we write \(\sigma_z\tau_z\ket{S}=-\ket{S}\), have we performed a measurement?}
\lecturerresponse{No. The equation records how an operator acts on a vector. It tells us that a measurement of the corresponding joint observable would be certain to return \(-1\), but no detector has been included merely by writing the eigenvalue equation.}
\end{classroomqa}

\begin{classroomqa}
\audiencequestion{Can an entangled pair be only partially correlated?}
\lecturerresponse{Yes. Perfect anticorrelation is a special property of the singlet for equal axes. Other entangled states can have weaker correlations, and even the singlet gives \(-\mathbf a\cdot\mathbf b\) rather than \(-1\) when the two axes differ. Entanglement concerns nonfactorization, not perfect agreement or disagreement in every chosen measurement.}
\end{classroomqa}

\section{Preparing the Singlet by Energy Relaxation}

The rotationally invariant two-spin operator
\begin{equation}
\boldsymbol{\sigma}\cdot\boldsymbol{\tau}
=\sigma_x\tau_x+\sigma_y\tau_y+\sigma_z\tau_z
\label{eq:l7-spin-dot-spin}
\end{equation}
organizes the four-dimensional state space. On the singlet,
\begin{equation}
(\boldsymbol{\sigma}\cdot\boldsymbol{\tau})\ket{S}
=-3\ket{S},
\label{eq:l7-singlet-dot-eigenvalue}
\end{equation}
while each of the three triplet states has eigenvalue \(+1\).

\begin{figure}[htbp]
\centering
\includegraphics[width=0.88\textwidth]{lecture_07_figure_05.png}
\caption{The rotationally invariant coupling \(\boldsymbol{\sigma}\cdot\boldsymbol{\tau}\) and its role in separating singlet and triplet sectors.}
\lectureframe{00:52:33}
\end{figure}

Suppose the interaction Hamiltonian contains
\begin{equation}
H_{\mathrm{int}}
=J\,\boldsymbol{\sigma}\cdot\boldsymbol{\tau}.
\label{eq:l7-exchange-hamiltonian}
\end{equation}
For antiferromagnetic coupling \(J>0\), the singlet lies below the triplet by
\begin{equation}
\Delta E=4J.
\end{equation}
Bring the spins together, allow the higher-energy components to decay by emitting a photon, a phonon, or some other excitation, and the pair relaxes toward the singlet. The spins may then be separated while preserving their joint state, provided the separation does not introduce uncontrolled spin rotations or close the energy gap too abruptly.

The same energy logic appears in atomic systems. In positronium, for example, the electron and positron spin coupling splits singlet and triplet configurations. An energy difference is also a mass difference through \(E=mc^2\), though spectroscopy usually reports the splitting as an energy or frequency.

\begin{classroomqa}
\audiencequestion{Is the singlet always the ground state of two interacting spins?}
\lecturerresponse{No. It is the ground state for the sign \(J>0\) in Equation~\eqref{eq:l7-exchange-hamiltonian}. Reversing the sign makes the triplet lower. The preparation argument depends on the actual interaction and its energy ordering, not merely on the existence of two spins.}
\end{classroomqa}

\begin{classroomqa}
\audiencequestion{Can \(\boldsymbol{\sigma}\cdot\boldsymbol{\tau}\) be measured by separately measuring all three components?}
\lecturerresponse{Not on one pair by three simultaneous local spin measurements: \(\sigma_x,\sigma_y,\sigma_z\) are incompatible, and so are the three \(\tau\) components. Its expectation value can be reconstructed from three subensembles measured along \(x,y,z\), or the singlet--triplet observable can be measured directly by a suitable joint apparatus.}
\end{classroomqa}

\begin{classroomqa}
\audiencequestion{After the spins separate, must they continue interacting to remain entangled?}
\lecturerresponse{No. Interaction is needed to prepare or alter the state, but once the singlet has been prepared, unitary evolution can preserve its entanglement without a continuing force between the particles. Local magnetic fields may rotate the two factors, yet local unitary rotations do not by themselves destroy entanglement.}
\end{classroomqa}

\section{The Reduced Density Matrix}

Return to Alice's local expectation value:
\begin{equation}
\expect{L_A}
=\sum_{a',a,b}
\Psi^*(a',b)L_{a'a}\Psi(a,b).
\end{equation}
Regroup the factors that do not belong to \(L\):
\begin{equation}
\expect{L_A}
=\sum_{a',a}
\left[
\sum_b\Psi(a,b)\Psi^*(a',b)
\right]L_{a'a}.
\label{eq:l7-regrouped-expectation}
\end{equation}
The expression in brackets is Alice's reduced density matrix,
\begin{equation}
(\rho_A)_{aa'}
=\sum_b\Psi(a,b)\Psi^*(a',b).
\label{eq:l7-reduced-density-components}
\end{equation}
In basis-independent notation,
\begin{equation}
\rho_A
=\operatorname{Tr}_B\ket{\Psi}\bra{\Psi},
\label{eq:l7-partial-trace}
\end{equation}
and every Alice-local expectation value is
\begin{equation}
\expect{L_A}
=\operatorname{Tr}_A(\rho_A L).
\label{eq:l7-density-expectation}
\end{equation}

\begin{figure}[htbp]
\centering
\includegraphics[width=0.88\textwidth]{lecture_07_figure_06.png}
\caption{Alice's density matrix obtained by summing the joint amplitudes over Bob's basis label.}
\lectureframe{01:02:33}
\end{figure}

The reduced density matrix is Hermitian, positive, and normalized:
\begin{equation}
\rho_A^\dagger=\rho_A,
\qquad
\bra{\chi}\rho_A\ket{\chi}\geq0,
\qquad
\operatorname{Tr}\rho_A=1.
\label{eq:l7-density-properties}
\end{equation}
It contains exactly the information needed for all measurements confined to Alice. If two global states produce the same \(\rho_A\), no Alice-only experiment can distinguish them.

\begin{classroomqa}
\audiencequestion{Does Alice measure the density matrix as one observable?}
\lecturerresponse{Not in a single shot. The density matrix is the object from which all local measurement probabilities are computed. Given many identically prepared systems, Alice can reconstruct it by tomography: she measures a sufficient set of observables and infers the matrix from their expectation values.}
\end{classroomqa}

\begin{classroomqa}
\audiencequestion{Why sum over Bob rather than select one of Bob's states?}
\lecturerresponse{Because Alice has no outcome information from Bob. The partial trace sums over a complete Bob basis and gives the unconditioned local state. Selecting one Bob outcome instead would produce a conditional state, which is a different operation and requires knowing that outcome.}
\end{classroomqa}

For a product state \(\Psi(a,b)=\psi_A(a)\phi_B(b)\),
\begin{align}
(\rho_A)_{aa'}
&=\sum_b
\psi_A(a)\phi_B(b)
\psi_A^*(a')\phi_B^*(b) \\
&=\psi_A(a)\psi_A^*(a')
\sum_b|\phi_B(b)|^2 \\
&=\psi_A(a)\psi_A^*(a').
\end{align}
Thus
\begin{equation}
\rho_A=\ket{\psi_A}\bra{\psi_A}.
\label{eq:l7-pure-projector}
\end{equation}
This is a rank-one projector. Acting on \(\ket{\psi_A}\) gives
\begin{equation}
\rho_A\ket{\psi_A}
=\ket{\psi_A}\braket{\psi_A}{\psi_A}
=\ket{\psi_A},
\label{eq:l7-density-on-state}
\end{equation}
while for every \(\ket{\chi}\) orthogonal to \(\ket{\psi_A}\),
\begin{equation}
\rho_A\ket{\chi}
=\ket{\psi_A}\braket{\psi_A}{\chi}
=0.
\label{eq:l7-density-on-orthogonal}
\end{equation}

\begin{figure}[htbp]
\centering
\includegraphics[width=0.88\textwidth]{lecture_07_figure_07.png}
\caption{A product-state density matrix acting as a projector: eigenvalue \(1\) along \(\ket{\psi_A}\) and \(0\) on orthogonal vectors.}
\lectureframe{01:09:43}
\end{figure}

\section{The Density-Matrix Spectrum and Entanglement}

The eigenvalues of \(\rho_A\) are nonnegative and sum to one. For a pure state of the whole:
\begin{itemize}
\item exactly one nonzero eigenvalue, equal to \(1\), means the state is a product;
\item more than one nonzero eigenvalue means the state is entangled;
\item equal eigenvalues across the full available local support give maximal entanglement.
\end{itemize}
The statement is most transparent from a Schmidt decomposition,
\begin{equation}
\begin{aligned}
\ket{\Psi}_{AB}
&=\sum_{r=1}^{R}\sqrt{\lambda_r}
\ket{r}_A\ket{r}_B,\\
\lambda_r&\geq0,
\qquad
\sum_r\lambda_r=1.
\end{aligned}
\label{eq:l7-schmidt}
\end{equation}
for which
\begin{equation}
\rho_A=\sum_{r=1}^{R}\lambda_r\ket{r}_A\bra{r}.
\label{eq:l7-schmidt-density}
\end{equation}
The Schmidt rank \(R\) is one for a product state and exceeds one for an entangled pure state.

For the singlet,
\begin{equation}
\rho_A
=\operatorname{Tr}_B\ket{S}\bra{S}
=\frac12\ket{u}\bra{u}
+\frac12\ket{d}\bra{d}
=\frac12 I.
\label{eq:l7-singlet-reduced-density}
\end{equation}
Its eigenvalues are
\begin{equation}
\operatorname{spec}(\rho_A)
=\left\{\frac12,\frac12\right\}.
\end{equation}
Alice's local state is maximally mixed even though the state of the pair is pure.

The spectrum suggests a quantitative measure. The von Neumann entropy
\begin{equation}
S(\rho_A)
=-\operatorname{Tr}(\rho_A\log\rho_A)
=-\sum_r\lambda_r\log\lambda_r
\label{eq:l7-entanglement-entropy}
\end{equation}
vanishes for a product pure state and reaches \(\log d_A\) for a maximally entangled state when \(d_A\leq d_B\). The detailed use of entropy can be deferred; the spectrum already supplies the essential distinction.

\begin{classroomqa}
\audiencequestion{Must every eigenvalue of an entangled reduced density matrix be nonzero?}
\lecturerresponse{No. More than one nonzero eigenvalue is enough. If Alice's space is large, an entangled state may occupy only a lower-dimensional Schmidt subspace, leaving additional zero eigenvalues. Maximal entanglement means equal eigenvalues over the entire smaller subsystem, not merely equality among an arbitrarily chosen subset.}
\end{classroomqa}

\begin{classroomqa}
\audiencequestion{Does a mixed reduced density matrix by itself prove that Alice is entangled with Bob?}
\lecturerresponse{Only if the global state is known to be pure. A mixed local state can also arise from an ordinary classical mixture in a mixed global state. For a pure bipartite state, however, local mixedness is exactly equivalent to entanglement across the bipartition.}
\end{classroomqa}

\begin{classroomqa}
\audiencequestion{Can all moments of Alice's observables determine her probability distributions?}
\lecturerresponse{In a finite-dimensional system, expectation values of an operator basis determine \(\rho_A\), and \(\rho_A\) determines every local outcome distribution through the Born rule. The reconstruction requires an ensemble of identically prepared copies; one isolated event cannot reveal a density matrix.}
\end{classroomqa}

\section{Measurement as an Entangling Interaction}

Replace Alice and Bob by a spin and a measuring apparatus. Let the apparatus have a ready state \(\ket{0}_D\) and a triggered state \(\ket{1}_D\). An ideal interaction that records whether the spin is up can be represented by
\begin{align}
\ket{u}\ket{0}_D&\longrightarrow\ket{u}\ket{1}_D,\\
\ket{d}\ket{0}_D&\longrightarrow\ket{d}\ket{0}_D.
\label{eq:l7-measurement-basis-map}
\end{align}
The detector must begin in a known ready state; otherwise its final state would not provide an unambiguous record.

Linearity fixes the evolution of a superposition:
\begin{equation}
\left(\alpha\ket{u}+\beta\ket{d}\right)\ket{0}_D
\longrightarrow
\alpha\ket{u}\ket{1}_D
+\beta\ket{d}\ket{0}_D.
\label{eq:l7-measurement-entanglement}
\end{equation}
Unless one coefficient vanishes, the final state is entangled. The apparatus has acquired information precisely because its pointer state is correlated with the spin.

\begin{figure}[htbp]
\centering
\includegraphics[width=0.88\textwidth]{lecture_07_figure_08.png}
\caption{An ideal measurement correlates spin alternatives with distinct pointer states and thereby creates entanglement.}
\lectureframe{01:20:20}
\end{figure}

If the spin alone is retained and the detector is ignored, its reduced state is
\begin{equation}
\rho_{\mathrm{spin}}
=|\alpha|^2\ket{u}\bra{u}
+|\beta|^2\ket{d}\bra{d},
\label{eq:l7-postmeasurement-reduced-state}
\end{equation}
assuming orthogonal pointer records. The interference terms have disappeared from the local description because they are stored in correlations with the apparatus.

There are now two complementary descriptions. If the detector outcome is read and one conditions on the record, the spin is assigned the corresponding branch: \(\ket{u}\) after a trigger or \(\ket{d}\) after no trigger. This is the collapse description. If detector and spin are treated as one closed system, Equation~\eqref{eq:l7-measurement-entanglement} is ordinary unitary evolution and no branch has been removed from the joint state.

\begin{classroomqa}
\audiencequestion{Is measurement something beyond entanglement?}
\lecturerresponse{At the level of system--apparatus dynamics, measurement is an interaction that establishes a reliable correlation. Collapse is the conditional description used after a definite record is selected. Whether one regards that update as fundamental or as relative to a chosen cut is an interpretive question; the operational predictions agree.}
\end{classroomqa}

\begin{classroomqa}
\audiencequestion{Why does an ideal measurement leave the spin eigenstate unchanged?}
\lecturerresponse{That is the repeatability requirement for this model. If the input is already an eigenstate of the measured observable, the apparatus should record its value without changing that value. Real measurements can be destructive, but then their interaction must be modeled differently.}
\end{classroomqa}

\section{The Observer Hierarchy}

Once the apparatus is admitted as a quantum system, there is no natural reason to stop there. Light from the pointer interacts with an eye; signals from the eye become correlated with a brain; an observer may be watched by another observer. Schr\"odinger's cat and Wigner's friend dramatize the same nested structure:
\begin{equation}
\begin{gathered}
\text{spin}
\longrightarrow
\text{apparatus}
\longrightarrow
\text{cat}\\
\longrightarrow
\text{Schr\"odinger}
\longrightarrow
\text{Wigner}.
\end{gathered}
\label{eq:l7-observer-chain}
\end{equation}
Inside any chosen boundary, one may describe the chain by unitary evolution and increasing entanglement. Outside the boundary, one uses the outcome-conditioned language of collapse.

The crucial consistency condition is that the boundary may be moved. If the cat records spin up, Schr\"odinger must later record a cat that saw spin up, and Wigner must record Schr\"odinger recording that same history. Quantum mechanics does not identify a unique sharp line at which an objective event must be declared. Cats, eyes, and brains are made of physical degrees of freedom; exempting one of them from quantum mechanics would require an additional law.

Feynman's attitude toward the resulting discomfort was pragmatic. The conceptual problem may be so large that it is difficult even to state whether there is one, while the calculational rules continue to predict every experiment to which they apply. The purpose of the formalism is not to force one interpretation, but to make the source of the discomfort precise: complete knowledge of a whole can coexist with the absence of any pure-state description of its parts.

\begin{classroomqa}
\audiencequestion{If I am the cat, do I have to draw the line at myself?}
\lecturerresponse{You may use a collapse description for the definite record available to you. Another observer who includes you in a larger isolated system may use an entangled-state description. The demand is not that everyone draw the same line, but that all descriptions give consistent records whenever the observers later compare results.}
\end{classroomqa}

\begin{classroomqa}
\audiencequestion{Where does the objective event happen: at the atom, the cat, Schr\"odinger, or Wigner?}
\lecturerresponse{The unitary formalism does not select one of those locations as a privileged boundary. Perhaps ordinary words such as \emph{happen} are being pushed beyond the context in which they are precise. The formal tools tell us how correlations propagate and what each observer can predict; deciding how to interpret the final boundary is a separate task.}
\end{classroomqa}

\subsection{Weak coupling and imperfect readout}

The ideal mapping in Equation~\eqref{eq:l7-measurement-basis-map} produces perfectly distinguishable records. A weaker interaction may leave the apparatus states only partially distinguishable:
\begin{equation}
\alpha\ket{u}\ket{D_u}
+\beta\ket{d}\ket{D_d},
\qquad
0<|\braket{D_d}{D_u}|<1.
\label{eq:l7-weak-measurement}
\end{equation}
Then the apparatus receives only partial information and the entanglement is incomplete. This is the setting of a weak or unsharp measurement.

That situation must be distinguished from a bad downstream readout. The detector may have become fully correlated with the spin while the eye, cable, amplifier, or brain fails to transmit the record cleanly. In the first case the system--detector interaction is weak; in the second, a good record exists at the detector but later information transfer is inefficient. Detector efficiency can involve either stage, so the physical coupling must be identified before assigning the cause.

\begin{classroomqa}
\audiencequestion{If poor vision makes the pointer uncertain, was the quantum measurement weak?}
\lecturerresponse{Not necessarily. The spin may have been strongly and completely correlated with the apparatus, while the apparatus-to-eye or eye-to-brain channel was poor. A genuinely weak measurement means that the initial system--apparatus coupling itself produced only partially distinguishable records.}
\end{classroomqa}

\section{Locality and No Signaling}

The phrase \emph{action at a distance} should be given an operational meaning. In ordinary physics, a distant action changes what can be observed here. If the Sun vanished and Earth's motion changed immediately, rather than after a light-travel time, one could use the change to transmit information faster than light. That would be genuine nonlocal signaling.

Entanglement does not permit such signaling. Let Bob perform a local quantum operation represented by Kraus operators \(K_\mu\), with
\begin{equation}
\sum_\mu K_\mu^\dagger K_\mu=I_B.
\end{equation}
If Alice does not learn Bob's outcome, the final joint state is
\begin{equation}
\rho'_{AB}
=\sum_\mu
(I_A\otimes K_\mu)\rho_{AB}
(I_A\otimes K_\mu^\dagger).
\end{equation}
Taking the partial trace over Bob gives
\begin{equation}
\rho'_A
=\operatorname{Tr}_B\rho'_{AB}
=\operatorname{Tr}_B\rho_{AB}
=\rho_A.
\label{eq:l7-no-signaling}
\end{equation}
Every unconditioned Alice-local probability is therefore unchanged by Bob's choice of a trace-preserving local operation.

Conditioned states are different. If Bob measures and obtains outcome \(\mu\), Alice may assign
\begin{equation}
\rho_{A|\mu}
=\frac{
\operatorname{Tr}_B\!\left[
(I_A\otimes K_\mu)\rho_{AB}
(I_A\otimes K_\mu^\dagger)
\right]
}{p_\mu}.
\label{eq:l7-conditioned-state}
\end{equation}
But Alice cannot know which \(\rho_{A|\mu}\) applies until Bob communicates \(\mu\) by an ordinary signal. Before the message arrives, she must average over outcomes:
\begin{equation}
\sum_\mu p_\mu\rho_{A|\mu}=\rho_A.
\label{eq:l7-conditioned-average}
\end{equation}

\begin{classroomqa}
\audiencequestion{When does Alice's density matrix change if Bob measures first?}
\lecturerresponse{Her unconditioned density matrix does not change. Once Bob's result reaches her, she may replace it with the corresponding conditional density matrix. That update reflects newly available information and cannot be used before the classical message arrives.}
\end{classroomqa}

\begin{classroomqa}
\audiencequestion{Does Bob's measurement physically collapse Alice's state at the same instant?}
\lecturerresponse{Different interpretations describe the conditional update differently, but all agree on the testable statement: Alice's local outcome probabilities remain unchanged until information from Bob becomes available. No controllable instantaneous signal is produced.}
\end{classroomqa}

Imagine many singlet pairs distributed in numbered rows. On one member Bob measures \(z\), on another \(x\), on another \(y\), and so forth. Alice's local results remain an unbiased sequence for every row:
\begin{equation}
P_A(+|\mathbf a)=P_A(-|\mathbf a)=\frac12.
\label{eq:l7-alice-marginal}
\end{equation}
Nothing in her notebook reveals which axes Bob chose or which outcomes he obtained.

\begin{figure}[htbp]
\centering
\includegraphics[width=0.88\textwidth]{lecture_07_figure_09.png}
\caption{Many independently prepared pairs allow different measurement axes to be assigned to different rows; Alice's marginal statistics remain unchanged.}
\lectureframe{01:49:00}
\end{figure}

After a light-speed-limited telegram arrives listing Bob's axis and result for each row, Alice can sort her records. Equal-axis singlet measurements will then be opposite pair by pair. The message reveals correlations already present in the joint statistics; it does not retroactively alter Alice's earlier local frequency distribution.

\begin{classroomqa}
\audiencequestion{Could Alice infer Bob's choice from sufficiently many moments of her own data?}
\lecturerresponse{No. All Alice-local moments are fixed by \(\rho_A\), and Equation~\eqref{eq:l7-no-signaling} leaves that matrix unchanged. More copies improve her estimate of the same local density matrix; they do not reveal Bob's basis or outcome.}
\end{classroomqa}

\begin{classroomqa}
\audiencequestion{What if Bob chooses a different axis on every member of the ensemble?}
\lecturerresponse{Alice still sees the same local distribution. Once Bob sends the row-by-row settings and results, the two parties can compare conditional correlations for each pair of axes. The setting dependence appears in joint statistics, not in Alice's marginal statistics.}
\end{classroomqa}

\section{Bell's Question as a Simulation Game}

The distinction between no signaling and Bell nonlocality becomes clearer as a problem in classical simulation. Begin with one spin,
\begin{equation}
\ket{\psi}=\alpha_u\ket{u}+\alpha_d\ket{d},
\qquad
|\alpha_u|^2+|\alpha_d|^2=1.
\end{equation}
A classical program can store the two complex amplitudes, evolve them with the Schr\"odinger equation, and accept a user-selected measurement direction. When the measure command is issued, the program computes the Born probabilities, uses a pseudorandom number generator to select an outcome, displays that outcome, and updates its stored state to the corresponding eigenstate. Repeating the process reproduces the statistics of one spin.

Randomness alone is therefore not the dividing line between classical and quantum mechanics. Classical computers can generate sequences that are operationally random enough for a simulation, even if their pseudorandom algorithms are deterministic underneath.

For two spins in a product state, two separated computers also suffice. Alice's computer stores and evolves \(\ket{\psi_A}\); Bob's stores and evolves \(\ket{\phi_B}\). Each responds only to its local setting. Because the joint probabilities factorize, no communication is needed after the computers separate.

An entangled pair changes the problem. One may bring the two machines together before the experiment, load correlated hidden data, synchronize their clocks, and even give them identical random seeds. After separation, however, each output must depend only on its local measurement setting and the shared information supplied in advance. Bell's theorem shows that no such local hidden-variable scheme reproduces all quantum correlations for freely chosen settings.

\begin{figure}[htbp]
\centering
\includegraphics[width=0.88\textwidth]{lecture_07_figure_10.png}
\caption{The separated-computer model: local programs may share preparation data, but after separation they cannot communicate measurement settings.}
\lectureframe{02:00:00}
\end{figure}

If wires remain between the computers so that settings and outcomes can be exchanged, the simulation can reproduce the correlations. But then the pair is no longer a genuinely separated local model; it functions as one connected computer. The obstruction is not classical computation in the abstract. A single global classical machine can simulate finite quantum systems by storing the complete joint wave function, although the required memory generally grows exponentially with the number of subsystems. Bell's obstruction concerns a decomposition into independently operating local machines.

\begin{classroomqa}
\audiencequestion{Why can the two computers not simply use the same random-number sequence?}
\lecturerresponse{Shared randomness is already allowed in a local hidden-variable model. It can predetermine correlated answers, but each answer must work for every measurement setting that might later be chosen on that side. Bell inequalities express the constraints obeyed by all such prearranged local strategies; the quantum correlations violate those constraints.}
\end{classroomqa}

\begin{classroomqa}
\audiencequestion{Do the random numbers have to depend on the measurement directions?}
\lecturerresponse{A local output may depend on its own local direction and the shared hidden data, but not on the distant direction chosen after separation. Allowing dependence on both settings requires communication, retrocausality, or a failure of the usual measurement-independence assumption. Bell's theorem makes these assumptions explicit rather than hiding them in the word \emph{random}.}
\end{classroomqa}

\begin{classroomqa}
\audiencequestion{Would two separated quantum computers succeed where the classical programs fail?}
\lecturerresponse{They can succeed if they share entangled quantum degrees of freedom. In that case the computers are not simulating entanglement with local classical records; they physically possess it. Their local outcomes still cannot transmit a controllable faster-than-light message.}
\end{classroomqa}

\begin{classroomqa}
\audiencequestion{Can one ordinary classical computer simulate the entire pair?}
\lecturerresponse{Yes. It can store the four joint amplitudes and compute all measurement probabilities in one place. For many particles this direct simulation becomes expensive because the Hilbert-space dimension grows exponentially, but there is no Bell contradiction when all information is globally available inside one machine.}
\end{classroomqa}

\begin{classroomqa}
\audiencequestion{Does Bell's theorem therefore prove ordinary action at a distance?}
\lecturerresponse{It proves that local hidden-variable explanations satisfying Bell's assumptions cannot reproduce the observed correlations. It does not provide Alice or Bob with a superluminal signaling channel. Bell nonlocality and operational faster-than-light communication are distinct claims.}
\end{classroomqa}

\section{What Can Be Learned Locally?}

Local and joint observables must not be conflated. Bob can measure \(\sigma_x\) and Alice can measure \(\tau_x\); by later comparing records they estimate
\begin{equation}
\expect{\sigma_x\tau_x}.
\end{equation}
Neither local apparatus alone measures that product. The rotational scalar
\begin{equation}
\boldsymbol{\sigma}\cdot\boldsymbol{\tau}
=\sigma_x\tau_x+\sigma_y\tau_y+\sigma_z\tau_z
\end{equation}
requires either different subensembles for the three terms or a new joint apparatus that resolves the singlet and triplet sectors. One cannot obtain all three incompatible component products from one pair using the original local spin meters.

\begin{classroomqa}
\audiencequestion{Are there observables of the whole that Alice's and Bob's original apparatuses cannot measure separately?}
\lecturerresponse{Yes. A single product such as \(\sigma_x\tau_x\) can be estimated by two local \(x\)-measurements followed by comparison. The sum \(\boldsymbol{\sigma}\cdot\boldsymbol{\tau}\) cannot be assembled from simultaneous \(x,y,z\) readings on one pair. It needs a genuinely joint singlet--triplet measurement or separate ensembles for the incompatible terms.}
\end{classroomqa}

Bob cannot infer entanglement from one event. With many identically prepared copies, however, he can reconstruct his reduced density matrix. For a qubit,
\begin{equation}
\rho_B
=\frac12\left(
I+\mathbf r\cdot\boldsymbol{\sigma}
\right),
\label{eq:l7-bloch-density}
\end{equation}
and measurements along several axes determine the Bloch vector \(\mathbf r\). If \(|\mathbf r|=1\), Bob's local state is pure; if \(|\mathbf r|<1\), it is mixed. In the special case where the global preparation is known to be the same pure state on every trial, local mixedness proves that Bob is entangled with something outside his subsystem.

For a product of two pure spins, Bob can search for the polarization direction \(\mathbf n=\mathbf r/|\mathbf r|\). Along that axis, repeated measurements return one definite answer. For the singlet, \(\mathbf r=0\), so every axis remains locally random. The search never finds a polarization direction.

\begin{classroomqa}
\audiencequestion{Can Bob discover entanglement without speaking to Alice?}
\lecturerresponse{From an ensemble he can discover that his own reduced state is mixed. If he also knows that every pair was prepared in the same global pure state, he can infer entanglement across the boundary. Without that purity assumption, his mixed state could instead come from a classical preparation mixture, and local data alone would not settle the distinction.}
\end{classroomqa}

\begin{classroomqa}
\audiencequestion{Why are many identically prepared systems necessary?}
\lecturerresponse{Each spin measurement yields only one outcome and generally disturbs that copy. Tomography therefore distributes incompatible measurements across an ensemble prepared by the same procedure. The density matrix and its spectrum are statistical properties inferred from those repeated trials, not labels revealed by one particle.}
\end{classroomqa}

The logic now closes where it began. Entanglement is not the vague sharing of a state. It is the failure of a pure joint state to factor, visible locally as a non-projector reduced density matrix and jointly as correlations that cannot be reproduced by separated classical bookkeeping. The whole may be perfectly known while each part remains mixed. That fact makes measurement possible, preserves no signaling, and supplies the structure that Bell's theorem will test.
